\subsection{Sprint 2}

Durante esta etapa do projeto acabei participando diretamente no frontend do
aplicativo mobile ao invés do backend. Tendo isso em mente, foquei mais no estudo
do framework específico ao invés das tecnologias do projeto como um todo.

Trabalhei junto com o AGES I Gustavo Pretto que me auxiliou bastante
especificamente na criação de componentes e desenvolvimento de telas. No final,
tanto por trabalhos ou provas de outras disciplinas e por dificuldades do time, a
entrega desta sprint foi muito incompleta.
\\

\fbox{%
    \parbox{0.95\linewidth}{%
        \setlength{\parindent}{15pt}
            \itshape    A sprint 2 foi imensamente difícil, o período coincidiu com diversas atividades
    acumuladas e em datas próximas de outras cadeiras. Não só isso, mas também o
    fato que minha reunião one-on-one foi onde tive que encarar a realidade que não
    estava contribuindo para o time.
    
        Eu estava perdido no ambiente do grupo, tanto quanto nos estudos das
    tecnologias novas que nunca havia visto. Tentei bastante virar a chave nesta sprint,
    mas acho que algumas ideias do time e mentalidade de alguns não foram de muita
    ajuda também. A técnica de pair-programming foi utilizada para servir de ajuda aos
    iniciantes e agilizar o processo, tendo em vista que a entrega na sprint 1 foi
    incompleta.
    
        Minha dupla funcionou bem, mas outras que talvez precisassem de mais ajuda
    acabavam sendo um pouco teimosos para aceitar ajuda. Quando terminávamos
    alguma task e não havia outra imediatamente para trabalharmos, nossa dupla
    acabava ficando ociosa devido a esse problema. Talvez com um pouco mais de
    insistência minha, e outros com menos teimosia, poderíamos ter feito uma entrega
    melhor para o cliente.
    }%
}